% ============================================================
% §4  Dynamiczne stałe fundamentalne
% ============================================================
\section{Dynamiczne sta\l{}e fundamentalne}\label{sec:stale}

W~standardowej fizyce $c$, $\hbar$ i~$G$ s\k{a} uniwersalnymi sta\l{}ymi.
W~TGP, gdzie przestrze\'n jest generowana przez materi\k{e}, te
,,sta\l{}e'' staj\k{a} si\k{e} \textbf{polami zale\.znymi od lokalnej
g\k{e}sto\'sci wygenerowanej przestrzeni $\Phi$}. Jedyn\k{a} prawdziwie
sta\l{}\k{a} wielko\'sci\k{a} jest d\l{}ugo\'s\'c Plancka $\lP$.

% ----------------------------------------------------------
\subsection{Uzasadnienie: sta\l{}e opisuj\k{a} propagacj\k{e}}\label{ssec:uzasadnienie}

$c$ okre\'sla pr\k{e}dko\'s\'c propagacji \'swiat\l{}a,
$\hbar$ --- skal\k{e} efekt\'ow kwantowych,
$G$ --- si\l{}\k{e} oddzia\l{}ywania grawitacyjnego.
Wszystkie trzy opisuj\k{a} w\l{}asno\'sci
\emph{przestrzeni, w~kt\'orej zachodzi propagacja}.

Skoro przestrze\'n w~TGP nie jest sta\l{}ym t\l{}em, lecz dynamicznym
polem $\Phi$, naturalne jest, \.ze w\l{}asno\'sci propagacji
--- a~wi\k{e}c warto\'sci $c$, $\hbar$, $G$ --- zale\.z\k{a} od
lokalnego stanu tego pola.

% ----------------------------------------------------------
\subsection{Pr\k{e}dko\'s\'c \'swiat\l{}a $c(\Phi)$}\label{ssec:c}

\begin{axiom}[Pr\k{e}dko\'s\'c \'swiat\l{}a jako pole]\label{ax:c}
Pr\k{e}dko\'s\'c \'swiat\l{}a jest lokalnym polem zale\.znym od $\Phi$:
\begin{equation}\label{eq:cPhi}
    c(\Phi) = c_0 \left(\frac{\PhiZero}{\Phi}\right)^{1/2},
\end{equation}
gdzie $c_0$ jest obserwowaną pr\k{e}dko\'sci\k{a} \'swiat\l{}a w~warunkach
referencyjnych ($\Phi = \PhiZero$).
\end{axiom}

Interpretacja:
\begin{itemize}[nosep]
  \item Im g\k{e}\'sciej wygenerowana przestrze\'n (wi\k{e}ksze $\Phi$),
        tym \textbf{wolniej} propaguje si\k{e} \'swiat\l{}o.
  \item Dla $\Phi \to \infty$: $c \to 0$ --- propagacja zanika
        (,,zamro\.zona'' przestrze\'n, wn\k{e}trze czarnej dziury).%
        \footnote{W~sformu\l{}owaniu kanonicznym (Formulacja~A) dynamiczne
        sta\l{}e przyjmuj\k{a} posta\'c $c(g)$, $\hbar(g)$, $G(g)$
        z~$g = \Phi/\PhiZero < 1$ blisko materii.
        Czarna dziura: $g \to 0$ (nie $\Phi \to \infty$).
        Brak punktu duchowego w~Formulacji~A.
        Patrz: \texttt{tgp\_companion.tex}.}
  \item Dla $\Phi \to 0$: $c \to \infty$ --- ale ten limit jest
        niefizyczny, bo $\Phi = 0$ oznacza brak przestrzeni.
\end{itemize}

G\k{e}sta przestrze\'n spowalnia \'swiat\l{}o --- analogicznie do
o\'srodka materialnego o~du\.zym wsp\'o\l{}czynniku za\l{}amania.
Ale tu nie ma o\'srodka: to sama \emph{tkanka przestrzeni}
jest g\k{e}\'sciejsza.

\begin{remark}[Interpretacja $c(\Phi)$: pr\k{e}dko\'s\'c
lokalna fizyczna]\label{rem:c-interpretacja}
W~standardowej OTW lokalna pr\k{e}dko\'s\'c \'swiat\l{}a
mierzona przyrz\k{a}dami obserwatora w~jego bezpo\'srednim
otoczeniu wynosi \emph{zawsze} $c_0$ (zasada r\'ownowa\.zno\'sci).
Zmienno\'s\'c $c$ jest widoczna dopiero w~por\'ownaniu
mi\k{e}dzy \textbf{r\'o\.znymi regionami}.

W~TGP aksjomat~\ref{ax:c} definiuje $c(\Phi)$ jako
\textbf{lokaln\k{a} fizyczn\k{a} pr\k{e}dko\'s\'c \'swiat\l{}a}
--- pr\k{e}dko\'s\'c wynikaj\k{a}c\k{a} z~metryki efektywnej
$g_{\mu\nu}^{\mathrm{eff}}(\Phi)$ w~danym punkcie:
\begin{equation}\label{eq:c-lok}
    c_{\mathrm{lok}}(\Phi)
    = c_0\sqrt{\frac{\PhiZero}{\Phi}}.
\end{equation}
Wielko\'s\'c ta r\'o\.zni si\k{e} od pr\k{e}dko\'sci
koordynatowej. Dla metryki eksponencjalnej
(stw.~\ref{prop:conformal-unique}):
$c_{\mathrm{koord}}^2 = c_0^2\,e^{-4\delta\Phi/\PhiZero}$,
natomiast $c_{\mathrm{lok}}^2 = c_0^2\,e^{-2\delta\Phi/\PhiZero}
\approx c_0^2\,\PhiZero/\Phi$.
To w\l{}a\'snie $c_{\mathrm{lok}}$ jest tre\'sci\k{a}
aksjomatu~\ref{ax:c}.

Obserwator \emph{operacyjnie lokalny} (zanurzony w~regionie o~$\Phi$)
u\.zywaj\k{a}c przyrz\k{a}d\'ow zbudowanych z~materii
podlegaj\k{a}cej temu samemu $\Phi$ --- mierzy $c_0$
(samokonsystencja), bo wszystkie jego standardy pomiarowe
skaluj\k{a} si\k{e} sp\'ojnie ($c$, $\hbar$, $G$ zale\.z\k{a}
od $\Phi$ w~spos\'ob zachowuj\k{a}cy $\lP$).
Zmienno\'s\'c $c(\Phi)$ jest zatem
\emph{relacyjna}: ujawnia si\k{e} tylko przy por\'ownywaniu
region\'ow o~r\'o\.znym $\Phi$.

To r\'o\.zni TGP od standardowej OTW, gdzie zmienna
pr\k{e}dko\'s\'c koordynatowa jest artefaktem wyboru
uk\l{}adu wsp\'o\l{}rz\k{e}dnych, natomiast w~TGP $c(\Phi)$ jest
\textbf{fizyczn\k{a}} konsekwencj\k{a} r\'o\.znej
g\k{e}sto\'sci wygenerowanej przestrzeni. R\'o\.znica
jest testowalna: w~TGP stosunek $c(\Phi_1)/c(\Phi_2)$
jest \textbf{obserwowalny} i~zale\.zy od $\Phi_1/\Phi_2$;
w~OTW analogiczny stosunek zale\.zy od \emph{krzywizny},
nie od \.zadnego pola skalarnego.
\end{remark}

% ----------------------------------------------------------
\subsection{Sta\l{}a Plancka $\hbar(\Phi)$}\label{ssec:hbar}

\begin{axiom}[Emergentna sta\l{}a Plancka]\label{ax:hbar}
Efektywna sta\l{}a Plancka jest polem:
\begin{equation}\label{eq:hbarPhi}
    \hbar(\Phi) = \hbar_0 \left(\frac{\PhiZero}{\Phi}\right)^{1/2},
\end{equation}
gdzie $\hbar_0$ jest obserwowaną sta\l{}\k{a} Plancka.
\end{axiom}

Interpretacja:
\begin{itemize}[nosep]
  \item W~g\k{e}stej przestrzeni ($\Phi \gg \PhiZero$):
        $\hbar \to 0$ --- efekty kwantowe s\k{a} t\l{}umione.
        Materia zachowuje si\k{e} klasycznie nawet na poziomie
        pojedynczych cz\k{a}stek.
  \item W~rzadkiej przestrzeni ($\Phi \sim \PhiZero$):
        $\hbar \sim \hbar_0$ --- standardowa mechanika kwantowa.
\end{itemize}

To daje \textbf{dwie drogi do granicy klasycznej}:
\begin{enumerate}[nosep]
  \item Droga statystyczna: $N \to \infty$ (wiele cz\k{a}stek,
        fluktuacje wzgl\k{e}dne zanikaj\k{a}).
  \item Droga przestrzenna: $\Phi \to \infty$ (g\k{e}sta
        przestrze\'n t\l{}umi kwantowo\'s\'c).
\end{enumerate}

\begin{remark}[Czarna dziura a~kwantowo\'s\'c]
Blisko czarnej dziury $\Phi$ jest ogromne, wi\k{e}c $\hbar \to 0$:
materia staje si\k{e} \textbf{klasyczna} w~silnym polu grawitacyjnym.
To jest przeciwie\'nstwo standardowego oczekiwania (w~OTW
efekty kwantowe powinny \emph{narastać} przy osobliwo\'sci).
W~TGP osobliwo\'s\'c klasyczna jest \emph{naturalnie usuni\k{e}ta}:
materia nie mo\.ze ,,skondensowa\'c si\k{e} w~punkt'', bo przestrze\'n
wok\'o\l{} niej jest zbyt g\k{e}sta.
\end{remark}

% ----------------------------------------------------------
\subsection{Sta\l{}a grawitacyjna $G(\Phi)$}\label{ssec:G}

\begin{axiom}[Dynamiczna sta\l{}a grawitacyjna]\label{ax:G}
Sta\l{}a grawitacyjna jest polem wyznaczonym przez warunek
sta\l{}o\'sci d\l{}ugo\'sci Plancka:
\begin{equation}\label{eq:GPhi}
    G(\Phi) = G_0 \cdot \frac{\PhiZero}{\Phi},
\end{equation}
gdzie $G_0$ jest obserwowaną sta\l{}\k{a} grawitacyjn\k{a}.
\end{axiom}

Interpretacja:
\begin{itemize}[nosep]
  \item W~g\k{e}stej przestrzeni: $G$ maleje --- grawitacja jest
        \textbf{s\l{}absza}. Wi\k{e}cej przestrzeni ,,rozcieńcza''
        oddzia\l{}ywanie.
  \item W~rzadkiej przestrzeni: $G$ ro\'snie --- grawitacja jest
        silniejsza (ale te\.z mniej przestrzeni).
\end{itemize}

% ----------------------------------------------------------
\subsection{Wyprowadzenie sta\l{}ych z~metryki i~substratu}\label{ssec:stalych-deriv}
% ----------------------------------------------------------

Aksjomaty~\ref{ax:c}--\ref{ax:G} powy\.zej s\k{a} sformu\l{}owane
jako aksjomaty dla wygody pedagogicznej --- ich postulowanie nie jest
konieczne. Poni\.zej pokazujemy, \.ze warto\'sci $c(\Phi)$, $\hbar(\Phi)$,
$G(\Phi)$ \textbf{wynikaj\k{a}} z~metryki pomostowej i~warunku
sta\l{}o\'sci d\l{}ugo\'sci Plancka. Status aksjomat\'ow A6--A8
nale\.zy traktowa\'c jako \textbf{wygod\k{e} prezentacyjn\k{a}},
a~ich tre\'s\'c jako Propozycje (\ref{prop:c-from-metric},
\ref{prop:G-from-lP}).

\begin{proposition}[$c(\Phi)$ z~metryki pomostowej]\label{prop:c-from-metric}
Metryka pomostowa TGP (tw.~\ref{thm:metric-from-substrate-full},
sek.~\ref{ssec:metric-from-budget}) jest \textbf{jednoznacznie}
wyznaczona przez warunek antypodyczny $f \cdot h = 1$,
PPN~$\gamma = \beta = 1$ oraz $\ell_P = \text{const}$:
\begin{equation}\label{eq:metric-bridge-c}
  ds^2 = -\frac{c_0^2}{\psi}\,dt^2
         + \psi\,\delta_{ij}\,dx^i dx^j,
  \qquad \psi \equiv \frac{\Phi}{\PhiZero}.
\end{equation}
Pr\k{e}dko\'s\'c koordynatowa \'swiat\l{}a ($ds^2 = 0$, ruch promieniowy):
\begin{equation}\label{eq:c-from-metric-deriv}
  c_{\rm koord}^2 = \frac{|g_{tt}|}{g_{rr}}
  = \frac{c_0^2/\psi}{\psi}
  = \frac{c_0^2}{\psi^2}
  \quad\Longrightarrow\quad
  c_{\rm koord} = \frac{c_0}{\psi}
  = c_0\,\frac{\PhiZero}{\Phi}.
\end{equation}
Pr\k{e}dko\'s\'c lokalna (odleg\l{}o\'s\'c w\l{}a\'sciwa / czas koordynatowy):
\begin{equation}\label{eq:c-lok-derived}
  c_{\rm lok} = \sqrt{h}\cdot\frac{dr}{dt}
  = \sqrt{\psi}\cdot\frac{c_0}{\psi}
  = \frac{c_0}{\sqrt{\psi}}
  = c_0\sqrt{\frac{\PhiZero}{\Phi}}.
\end{equation}
Wynik~\eqref{eq:c-lok-derived} jest \textbf{dok\l{}adnie} r\'owny
aksjomat~\ref{ax:c} ($c(\Phi) = c_0\sqrt{\PhiZero/\Phi}$),
kt\'ory definiuje $c_{\rm lok}$, nie $c_{\rm koord}$.
Aksjomat jest zatem \textbf{konsekwencj\k{a}} metryki pomostowej,
nie niezale\.znym postulatem.

\medskip\noindent
\textbf{Trzy pr\k{e}dko\'sci \'swiat\l{}a w~TGP:}
\begin{center}\small
\begin{tabular}{@{}lll@{}}
\toprule
\textbf{Symbol} & \textbf{Warto\'s\'c} & \textbf{Interpretacja} \\
\midrule
$c_{\rm proper}$ & $c_0$ (zawsze)
  & lokalna, obie miary w\l{}a\'sciwe \\
$c_{\rm lok}$ & $c_0/\sqrt{\psi}$
  & odl.~w\l{}a\'sciwa / czas koordynatowy (A6) \\
$c_{\rm koord}$ & $c_0/\psi$
  & koordynatowa ($= c_{\rm lok}^2/c_0$) \\
\bottomrule
\end{tabular}
\end{center}
\noindent
Fizyka nie zale\.zy od wyboru miary: $\gamma_{\rm PPN} = \beta_{\rm PPN} = 1$
w~obu przypadkach (postaci pot\k{e}gowej i~eksponencjalnej).

\medskip\noindent
\textbf{R\'ownowa\.zna posta\'c eksponencjalna} (do pierwszego rz\k{e}du
w~$\delta\Phi = \Phi - \PhiZero$, stw.~\ref{hyp:metric-exp}):
\begin{equation}
  ds^2 \approx -e^{-\delta\Phi/\PhiZero}\,c_0^2\,dt^2
         + e^{+\delta\Phi/\PhiZero}\,\delta_{ij}\,dx^i dx^j,
\end{equation}
kt\'ora daje $c_{\rm koord} = c_0\,e^{-\delta\Phi/\PhiZero}
\approx c_0\,(1 - \delta\Phi/\PhiZero) = c_0\sqrt{\PhiZero/\Phi}
+ \mathcal{O}((\delta\Phi/\PhiZero)^2)$.
Obie postaci (pot\k{e}gowa i~eksponencjalna) s\k{a} sp\'ojne
w~re\.zimie s\l{}abego pola; r\'o\.zni\k{a} si\k{e} dopiero
w~drugim rz\k{e}dzie post-newtonowskim, gdzie eksponencjalna
daje poprawne $\beta_{\rm PPN} = 1$
(uw.~\ref{rem:power-vs-exp}).
\end{proposition}

\begin{proposition}[$G(\Phi)$ z~warunku sta\l{}o\'sci $\ell_P$]\label{prop:G-from-lP}
Za\l{}\'o\.zmy:
\begin{enumerate}[nosep]
  \item $c(\Phi) = c_0(\PhiZero/\Phi)^{1/2}$ (wynika z~metryki, prop.~\ref{prop:c-from-metric}),
  \item $\hbar(\Phi) = \hbar_0(\PhiZero/\Phi)^{b}$ (ogolna pot\k{e}gowa),
  \item $\ell_P^2 = \hbar G/c^3 = \mathrm{const}$.
\end{enumerate}
Warunek (3) daje: $G(\Phi) \propto c^3/\hbar \propto (\PhiZero/\Phi)^{3/2 - b}/(\PhiZero/\Phi)^b
= (\PhiZero/\Phi)^{3/2 - 2b}$.
Aby $c/\hbar = c_0/\hbar_0 = \mathrm{const}$ (masa cz\k{a}stki niezale\.zna od $\Phi$,
warunek~W3 tw.~\ref{thm:exponents}): $b = 1/2$. Wtedy:
\begin{equation}
  G(\Phi) = G_0\,\frac{\PhiZero}{\Phi},
\end{equation}
czyli aksjomat~\ref{ax:G} wynika z~$\ell_P = \mathrm{const}$ + metryki + sp\'ojno\'sci.
Status aksjomat\'ow A6--A8: \textbf{Propozycje wyprowadzone z~A1, A3, A6}.
\end{proposition}

\begin{remark}[Status epistemiczny stałych]\label{rem:stalych-status}
W~\'swietle propozycji~\ref{prop:c-from-metric}--\ref{prop:G-from-lP}
i~twierdzenia~\ref{thm:exponents}:
\begin{center}\small
\begin{tabular}{@{}llll@{}}
\toprule
\textbf{Stała} & \textbf{Forma} & \textbf{Formalny status} & \textbf{Źródło} \\
\midrule
$c(\Phi)$ & $c_0(\Phi_0/\Phi)^{1/2}$ & Propozycja &
  metryka pomostowa (stw.~\ref{prop:conformal-unique}) \\
$\hbar(\Phi)$ & $\hbar_0(\Phi_0/\Phi)^{1/2}$ & Propozycja &
  $c/\hbar = \mathrm{const}$ (warunek W3) \\
$G(\Phi)$ & $G_0\,\Phi_0/\Phi$ & Twierdzenie &
  $\ell_P = \mathrm{const}$ (tw.~\ref{thm:lP}) \\
$\ell_P$ & $\sqrt{\hbar_0 G_0/c_0^3}$ & Twierdzenie &
  tw.~\ref{thm:lP}: inwariant substratowy \\
\bottomrule
\end{tabular}
\end{center}
Wszystkie trzy ``sta\l{}e'' s\k{a} \textbf{wyprowadzone} z~jednego pola $\Phi$
i~warunków wewn\k{e}trznych TGP. Nie ma potrzeby postulowania ich
jako niezale\.znych aksjomat\'ow.
\end{remark}

% ----------------------------------------------------------
\subsection{Hierarchia uzasadnienia wyk\l{}adnik\'ow}\label{ssec:hierarchia-uzasad}

Zanim przejdziemy do formalnego dowodu jednoznaczno\'sci
wyk\l{}adnik\'ow, warto wskaza\'c \textbf{hierarchi\k{e} logiczn\k{a}}
prowadz\k{a}c\k{a} od ontologii TGP do konkretnych pot\k{e}g
skalowania:

\begin{remark}[Trzy kroki do $(a,b,g) = (1/2, 1/2, 1)$]\label{rem:hierarchia-skalowania}
\begin{enumerate}
  \item \textbf{Metryka efektywna z~$\Phi$}
    (stw.~\ref{prop:conformal-unique}):
    Aksjomat~\ref{ax:przestrzen} (jedynym polem jest $\Phi$)
    + warunek Minkowskiego przy $\Phi = \PhiZero$
    + interpretacja fizyczna pr\k{e}dko\'sci \'swiat\l{}a
    jednoznacznie daj\k{a} antypodaln\k{a} metryk\k{e}
    z~wyk\l{}adnikiem $p = 1/2$. St\k{a}d:
    $c(\Phi) = c_0(\PhiZero/\Phi)^{1/2}$, czyli $a = 1/2$.

  \item \textbf{Wymuszenie sta\l{}o\'sci $\lP$}
    (tw.~\ref{thm:lP}):
    D\l{}ugo\'s\'c Plancka jest jedyn\k{a} fundamentaln\k{a}
    skal\k{a} TGP (niezale\.zn\k{a} od $\Phi$).
    Warunek $\lP^2 = \hbar G / c^3 = \mathrm{const}$
    daje zwi\k{a}zek $b + g - 3a = 0$.

  \item \textbf{Sp\'ojno\'s\'c propagatora}:
    Kwanty pola $\Phi$ (fonony przestrzenno\'sci) musz\k{a}
    mie\'c standardow\k{a} dyspersj\k{e} $\omega = ck$
    w~ramce lokalnej. Wymaga to, by d\l{}ugo\'s\'c fali
    de Broglie'a $\lambda_{\mathrm{dB}} = \hbar/(mc)$
    skalowa\l{}a si\k{e} tak samo jak $c^{-1}$, co daje $a = b$.
\end{enumerate}

Razem: $a = 1/2$ (krok~1), $b = a = 1/2$ (krok~3),
$g = 3a - b = 1$ (krok~2). \.Zaden wyk\l{}adnik nie jest
wstawiany r\k{e}cznie --- ka\.zdy wynika z~wewn\k{e}trznych
zasad TGP. Dowodzimy tego formalnie poni\.zej.
\end{remark}

% ----------------------------------------------------------
\subsection{Wyprowadzenie wyk\l{}adnik\'ow z~zasad
  wewn\k{e}trznych}\label{ssec:wykladniki}

Aksjomaty~\ref{ax:c}--\ref{ax:G} postuluj\k{a} konkretne
zale\.zno\'sci pot\k{e}gowe: $c \propto \Phi^{-1/2}$,
$\hbar \propto \Phi^{-1/2}$, $G \propto \Phi^{-1}$.
Poni\.zej pokazujemy, \.ze te wyk\l{}adniki s\k{a}
\textbf{jedynymi} zgodnymi z~trzema wewn\k{e}trznymi
wymaganiami TGP:

\begin{theorem}[Jednoznaczno\'s\'c wyk\l{}adnik\'ow]\label{thm:exponents}
Niech $c(\Phi) = c_0(\PhiZero/\Phi)^a$,
$\hbar(\Phi) = \hbar_0(\PhiZero/\Phi)^b$,
$G(\Phi) = G_0(\PhiZero/\Phi)^g$. Za\l{}\'o\.zmy:
\begin{enumerate}[nosep]
  \item[(W1)] \textbf{Sta\l{}o\'s\'c $\lP$:}
    $\lP^2 = \hbar G/c^3 = \mathrm{const}$
    (fundamentalna skala nie zale\.zy od $\Phi$).
  \item[(W2)] \textbf{Pr\k{e}dko\'s\'c \'swiat\l{}a z~metryki:}
    Metryka antypodaln\k{a} $f = (\PhiZero/\Phi)^p$,
    $h = (\Phi/\PhiZero)^p$ (stwierdzenie~\ref{prop:conformal-unique})
    daje $c^2 = c_0^2\,f/h
    = c_0^2\,(\PhiZero/\Phi)^{2p}$; aksjomat~\ref{ax:c}
    wymaga $2p = 1$, czyli $p = 1/2$, sk\k{a}d $a = 1/2$.
    (Dla metryki minimalnej $c_{\mathrm{koord}} = c_{\mathrm{lok}}$;
    dla eksponencjalnej por.\ stw.~\ref{prop:conformal-unique}.)
  \item[(W3)] \textbf{Sp\'ojno\'s\'c propagatora (uniwersalno\'s\'c
    fizyki lokalnej):}
    Propagator Kleina--Gordona cz\k{a}stki o~masie spoczynkowej~$m$
    ma w~ramce lokalnej posta\'c
    \[
      G(k) \;=\; \frac{1}{k^2 + m_{\mathrm{eff}}^2},
      \qquad
      m_{\mathrm{eff}} \;=\; \frac{mc(\Phi)}{\hbar(\Phi)},
    \]
    gdzie $m_{\mathrm{eff}}$ jest odwrotn\k{a} d\l{}ugo\'sci\k{a}
    Comptona --- fizycznym parametrem masowym
    (biegunem propagatora). Skalowanie:
    \[
      \frac{c(\Phi)}{\hbar(\Phi)}
      \;=\; \frac{c_0}{\hbar_0}
            \left(\frac{\PhiZero}{\Phi}\right)^{\!a-b}.
    \]
    Wymag\-amy, aby fizyka lokalna by\l{}a \textbf{uniwersalna}:
    masa mierzona lokalnie nie zale\.zy od po\l{}o\.zenia
    (tzn.\ $m_{\mathrm{eff}}$ jest sta\l{}e przy sta\l{}ym~$m$).
    St\k{a}d $a - b = 0$, czyli $a = b$.
\end{enumerate}
Wtedy $a = b = 1/2$ (z~W2 i~W3). Warunek W1 daje:
\begin{equation}
  b + g - 3a = 0
  \quad\Longrightarrow\quad
  g = 3a - b = 3\cdot\tfrac{1}{2} - \tfrac{1}{2} = 1.
\end{equation}
\textbf{Jedynym} rozwi\k{a}zaniem jest $(a,b,g) = (1/2,\;1/2,\;1)$
--- dok\l{}adnie aksjomaty~\ref{ax:c}--\ref{ax:G}.
\end{theorem}

\begin{proof}
\textbf{(W1)} Kombinacja $\lP^2 = \hbar G/c^3$ daje:
\[
  \lP^2 = \frac{\hbar_0\,G_0}{c_0^3}
  \left(\frac{\PhiZero}{\Phi}\right)^{b+g-3a}.
\]
Sta\l{}o\'s\'c $\lP$ wymaga $b + g - 3a = 0$.

\textbf{(W2)} Metryka konforemna z~stwierdzenia~\ref{prop:conformal-unique}
daje $a = 1/2$.

\textbf{(W3)} Stosunek $c(\Phi)/\hbar(\Phi) = (c_0/\hbar_0)(\PhiZero/\Phi)^{a-b}$.
Uniwersalno\'s\'c fizyki lokalnej (biegun propagatora $m_{\mathrm{eff}} = mc/\hbar$
niezale\.zny od $\Phi$) wymaga $a - b = 0$, czyli $b = a = 1/2$.

Podstawiaj\k{a}c do~(W1): $g = 3a - b = 3\cdot\frac{1}{2} - \frac{1}{2} = 1$.

Uk\l{}ad trzech r\'owna\'n jest liniowy i~niezdegenerowany
(macierz wsp\'o\l{}czynnik\'ow ma rz\k{a}d~3), wi\k{e}c
rozwi\k{a}zanie jest \textbf{jednoznaczne}.
\end{proof}

\begin{remark}[Konsekwencja: aksjomaty A6--A8 nie s\k{a} niezale\.zne]
Twierdzenie~\ref{thm:exponents} pokazuje, \.ze spo\'sr\'od
trzech aksjomat\'ow A6--A8 wystarczy postulowa\'c \textbf{jeden}
(np.~A6: $c \propto \Phi^{-1/2}$) wraz z~warunkami W1 i~W3.
Pozosta\l{}e dwa s\k{a} \emph{konsekwencjami}. Redukuje to
liczb\k{e} niezale\.znych za\l{}o\.ze\'n TGP.
\end{remark}

\begin{remark}[Jednoznaczno\'s\'c w~klasie pot\k{e}gowej]\label{rem:power-law-class}
Twierdzenie~\ref{thm:exponents} zak\l{}ada pot\k{e}gowe postacie
$c(\Phi)\propto\Phi^{-a}$, $\hbar(\Phi)\propto\Phi^{-b}$,
$G(\Phi)\propto\Phi^{-g}$. Jednoznaczno\'s\'c $(a,b,g)=(1/2,1/2,1)$
obowi\k{a}zuje \textbf{wewn\k{a}trz tej klasy}.
Rozwa\.zmy og\'olniejsze postacie, np.~logarytmiczne
$c = c_0 f(\ln\Phi/\PhiZero)$ lub saturuj\k{a}ce
$c = c_0/(1 + \Phi/\Phi_*)^a$. Warunki W1--W3 generuj\k{a} wtedy
uk\l{}ad r\'owna\'n funkcyjnych, kt\'orego jedynymi
\emph{ci\k{a}g\l{}ymi i~regularnym} rozwi\k{a}zaniami z~granicami
$c \to c_0$ dla $\Phi\to\PhiZero$ s\k{a} w\l{}a\'snie postacie
pot\k{e}gowe --- co wynika z~liniowo\'sci warunku $b+g-3a=0$ oraz
multiplikatywnej struktury skalowa\'n. Ma\l{}e odchylenia
$\delta f(\Phi) = \varepsilon\,h(\Phi)$ od klasy pot\k{e}gowej
prowadz\k{a} do zmian obserwabli bezwymiarowych rz\k{e}du
$O(\varepsilon\,\Phi_*/\PhiZero)$, co dla $\Phi_*\gg\PhiZero$
(realistyczne warunki) mo\.ze by\'c dost\k{e}pne testom
zmienno\'sci sta\l{}ych~\cite{Uzan2011}.
\end{remark}

% ----------------------------------------------------------
\subsection{Sta\l{}o\'s\'c d\l{}ugo\'sci Plancka}\label{ssec:lP}\label{ssec:naturalunits}

D\l{}ugo\'s\'c Plancka jest zdefiniowana jako:
\begin{equation}
    \lP^2 = \frac{\hbar\,G}{c^3}.
\end{equation}

Podstawiaj\k{a}c zale\.zno\'sci~\eqref{eq:cPhi}, \eqref{eq:hbarPhi},
\eqref{eq:GPhi}:
\begin{equation}
    \lP^2
    = \frac{\hbar_0\left(\frac{\PhiZero}{\Phi}\right)^{1/2}
            \cdot G_0\left(\frac{\PhiZero}{\Phi}\right)}
           {\left[c_0\left(\frac{\PhiZero}{\Phi}\right)^{1/2}\right]^3}
    = \frac{\hbar_0\,G_0}{c_0^3}
      \cdot \left(\frac{\PhiZero}{\Phi}\right)^{1/2+1-3/2}
    = \frac{\hbar_0\,G_0}{c_0^3}
      \cdot \left(\frac{\PhiZero}{\Phi}\right)^{0}
    = \frac{\hbar_0\,G_0}{c_0^3}.
\end{equation}

\begin{theorem}[Sta\l{}o\'s\'c skali Plancka]\label{thm:lP}
D\l{}ugo\'s\'c Plancka jest \textbf{niezmienna}:
\begin{equation}\label{eq:lP-const}
    \boxed{\lP = \sqrt{\frac{\hbar_0\,G_0}{c_0^3}} = \mathrm{const},}
\end{equation}
niezale\.znie od lokalnej warto\'sci $\Phi$. Jest to jedyna
prawdziwie fundamentalna skala TGP.
\end{theorem}

To jest \textbf{warunek sp\'ojno\'sci wewn\k{e}trznej}: pomimo
dynamiczno\'sci $c$, $\hbar$, $G$, skala kwantowej grawitacji
pozostaje sta\l{}a. \'Swiadczy to o~tym, \.ze zale\.zno\'sci
$c(\Phi)$, $\hbar(\Phi)$, $G(\Phi)$ nie s\k{a} niezale\.znymi
za\l{}o\.zeniami, lecz \textbf{jednym sp\'ojnym warunkiem}
wynikaj\k{a}cym z~natury pola $\Phi$.
W~istocie, twierdzenie~\ref{thm:exponents} pokazuje, \.ze
wyk\l{}adniki $(1/2,\;1/2,\;1)$ s\k{a} \textbf{jedynym}
rozwi\k{a}zaniem trzech warunk\'ow wewn\k{e}trznych TGP
(sta\l{}o\'s\'c $\lP$, metryka konforemna, sp\'ojno\'s\'c
propagatora).

% ----------------------------------------------------------
\subsection{Konsekwencje}\label{ssec:stale-kons}

\begin{enumerate}[nosep]
  \item \textbf{Nie ma problemu sta\l{}ej kosmologicznej}
        w~standardowym sensie: $G$ zmienia si\k{e} z~$\Phi$,
        wi\k{e}c ,,st a\l{}a'' kosmologiczna wyra\.zona przez $G$
        nie jest sta\l{}\k{a}.
  \item \textbf{Soczewkowanie grawitacyjne}: w~regionach
        o~wy\.zszym $\Phi$ \'swiat\l{}o porusza si\k{e} wolniej ---
        to daje soczewkowanie bez potrzeby krzywizny
        jako osobnego poj\k{e}cia.
  \item \textbf{Czerwone przesuni\k{e}cie}: foton id\k{a}cy
        z~regionu o~wy\.zszym $\Phi$ do ni\.zszego przyspiesza
        ($c$ ro\'snie), co powoduje przesuni\k{e}cie ku
        czerwieni --- odtwarza grawitacyjne czerwone
        przesuni\k{e}cie OTW.
  \item \textbf{Granica klasyczna} przy silnych polach:
        blisko masywnych obiekt\'ow $\hbar \to 0$, usuwaj\k{a}c
        potrzeb\k{e} kwantowej grawitacji w~silnym polu.
\end{enumerate}

% ----------------------------------------------------------
\subsection{Mapa obserwabilno\'sci dynamicznych sta\l{}ych}%
\label{ssec:obserwabilnosc-stalych}
% ----------------------------------------------------------

Zmiany $c(\Phi)$, $\hbar(\Phi)$, $G(\Phi)$ s\k{a} \textbf{relacyjne}:
lokalny obserwator zawsze mierzy warto\'sci referencyjne $c_0$, $\hbar_0$,
$G_0$, poniewa\.z r\'ownie\.z jego \textbf{wzorce} (linijki, zegary,
spektrometry) skaluj\k{a} si\k{e} z~$\Phi$ w~identyczny spos\'ob.
Efekty s\k{a} wykrywalne tylko przez \emph{por\'ownanie} dw\'och
region\'ow o~r\'o\.znym $\Phi$.

\begin{center}
\small
\begin{tabular}{@{}p{3.2cm}p{2.8cm}p{2.5cm}p{5.5cm}@{}}
\toprule
\textbf{Wielko\'s\'c} & \textbf{Zale\.zno\'s\'c od~$\Phi$}
  & \textbf{Lokalna obserwowalno\'s\'c}
  & \textbf{Jak wykry\'c} \\
\midrule
Pr\k{e}dko\'s\'c \'swiat\l{}a $c$
  & $c_0\sqrt{\PhiZero/\Phi}$
  & NIE (lokalnie zawsze $c_0$)
  & Por\'ownanie cz\k{e}sto\'sci foton\'ow z~region\'ow
    o~r\'o\.znym $\Phi$ (czerwone przesuni\k{e}cie) \\
\midrule
Sta\l{}a Plancka $\hbar$
  & $\hbar_0\sqrt{\PhiZero/\Phi}$
  & NIE
  & Zmiana widma atomowego mi\k{e}dzy regionami;
    poziom K6 (K6 OTWARTY) \\
\midrule
Sta\l{}a grawitacyjna $G$
  & $G_0\,\PhiZero/\Phi$
  & NIE (lokalnie $G_0$)
  & $G(z)$ z~BBN/CMB; SPARC galactic scaling; K5 ZAMKNI\k{E}TY \\
\midrule
D\l{}ugo\'s\'c Plancka $\ell_P$
  & $\sqrt{\hbar_0 G_0/c_0^3}$ --- \textbf{CONST}
  & TAK (inwariant)
  & Fundamentalna skala TGP; K7 ZAMKNI\k{E}TY \\
\midrule
Stosunek $G/c^3$
  & $G_0/c_0^3 \cdot (\PhiZero/\Phi)^{1/2}$
  & NIE (lokalnie const)
  & Amplituda fal grawitacyjnych (K20); LISA 2034+ \\
\midrule
Hubble $H(z)$
  & $H_{\Lambda\mathrm{CDM}}(z)/\sqrt{\phi(z)}$
  & TAK (relacja $d_L$--$z$)
  & DESI, Euclid; K2 OTWARTY \\
\midrule
Masa substratu $m_{\mathrm{sp}}$
  & $\sqrt\gamma \cdot \PhiZero$ (sta\l{}a)
  & TAK (wyznacza skal\k{e} Higgsa)
  & K15 ZAMKNI\k{E}TY (12/12 PASS) \\
\bottomrule
\end{tabular}
\end{center}

\begin{remark}[Zasada obserwabilno\'sci relacyjnej]\label{rem:relacyjna}
  W~TGP \emph{obserwowalne s\k{a} tylko stosunki wielko\'sci z~r\'o\.znych
  region\'ow przestrzeni}.
  Nie istnieje obiektywna ``absolutna'' warto\'s\'c $c$, $\hbar$ ani $G$
  --- s\k{a} one numerycznie r\'owne referencyjnym $c_0$, $\hbar_0$, $G_0$
  \emph{z~definicji} dla obserwatora lokalnego. Efekty TGP objawiaj\k{a}
  si\k{e} jako anomalie w~przeskalowaniu standardowych relacji
  (Hubble, wielko\'s\'c kszta\l{}t galaktyk, czas \.zycia gwiazd neutronowych).
\end{remark}

% ----------------------------------------------------------
\subsection{Termodynamiczna sp\'ojno\'s\'c dynamicznych sta\l{}ych}%
\label{ssec:thermo-consistency}

Dynamiczne $c(\Phi)$, $\hbar(\Phi)$, $G(\Phi)$ zmieniaj\k{a}
wielko\'sci termodynamiczne --- np.\ sta\l{}\k{a}
Stefana--Boltzmanna $\sigma_{\mathrm{SB}}$ zale\.zy od $\hbar$ i~$c$.
Poni\.zej pokazujemy, \.ze to\.zsamo\'sci termodynamiczne
pozostaj\k{a} spe\l{}nione \textbf{lokalnie}, a~zmienione
warto\'sci globalnych wielko\'sci s\k{a} \emph{predykcjami}, nie
nisp\'ojno\'sciami.

\begin{proposition}[Lokalna sp\'ojno\'s\'c termodynamiczna]%
\label{prop:thermo-consistency}
Niech obserwator znajduje si\k{e} w~regionie o~sta\l{}ym $\Phi$.
Wszystkie jego przyrz\k{a}dy pomiarowe podlegaj\k{a} temu samemu
$\Phi$ (aksjomat~\ref{ax:metric-coupling}). Wtedy:
\begin{enumerate}[nosep]
  \item Stosunek $k_B T / (\hbar\,\omega)$ jest niezmienniczy
        wzgl\k{e}dem $\Phi$-przeskalowania, gdy cz\k{e}sto\'s\'c
        $\omega \propto c/\lambda \propto \Phi^{-1/2}$
        i~temperatura lokalna dostosowuje si\k{e} jak
        $T_{\mathrm{lok}} \propto \hbar\,\omega / k_B$.
  \item Czynniki Boltzmanna
        $\exp(-E/(k_B T))$ s\k{a} zachowane lokalnie,
        bo $E \propto \hbar\,\omega$ skaluje si\k{e} sp\'ojnie z~$T$.
  \item R\'ownowaga termodynamiczna (rozk\l{}ad Plancka,
        rozk\l{}ad Maxwella--Boltzmanna) ma \textbf{identyczn\k{a}}
        posta\'c funkcyjn\k{a} dla ka\.zdego $\Phi$.
\end{enumerate}
Obserwator lokalny mierzy te same prawa termodynamiki niezale\.znie
od $\Phi$ --- zmienno\'s\'c jest \emph{relacyjna}, tak jak
w~przypadku $c(\Phi)$ (por.\ uwaga~\ref{rem:c-interpretacja}).
\end{proposition}

\begin{proof}
Sta\l{}a Boltzmanna $k_B$ jest wielko\'sci\k{a} statystyczn\k{a}
(przelicznik temperatura--energia) i~\textbf{nie zale\.zy} od
$\Phi$. Energia fotonu $E = \hbar(\Phi)\,\omega(\Phi)$, gdzie
$\omega = 2\pi c(\Phi)/\lambda$. Poniewa\.z $\hbar \propto
\Phi^{-1/2}$ i~$c \propto \Phi^{-1/2}$, mamy
$E \propto \Phi^{-1}$. Temperatura lokalna,
wyznaczona przez $k_B T = E_{\mathrm{char}}$, skaluje si\k{e}
tak samo: $T_{\mathrm{lok}} \propto \Phi^{-1}$.
St\k{a}d stosunek $E/(k_B T)$ jest niezale\.zny od $\Phi$,
a~wszystkie rozk\l{}ady r\'ownowagowe zachowuj\k{a} posta\'c.
\end{proof}

\begin{remark}[Skalowanie Stefana--Boltzmanna jako predykcja]%
\label{rem:stefan-boltzmann}
Sta\l{}a Stefana--Boltzmanna $\sigma_{\mathrm{SB}}
= \pi^2 k_B^4 / (60\,\hbar^3 c^2)$ skaluje si\k{e} z~$\Phi$:
\begin{equation}\label{eq:sigma-SB-Phi}
  \sigma_{\mathrm{SB}}(\Phi)
  = \sigma_{\mathrm{SB},0}
    \left(\frac{\Phi}{\PhiZero}\right)^{5/2},
\end{equation}
bo $\hbar^{-3} \propto \Phi^{3/2}$ i~$c^{-2} \propto \Phi$.
Nie jest to nisp\'ojno\'s\'c, lecz \textbf{predykcja TGP}:
g\k{e}sto\'s\'c energii promieniowania cieplnego jest
\emph{wi\k{e}ksza} w~regionach o~wy\.zszym $\Phi$
(g\k{e}\'sciejsza przestrze\'n $\to$ wi\k{e}cej stopni swobody
na jednostk\k{e} obj\k{e}to\'sci koordynatowej).
Lokalny obserwator tego nie zauwa\.zy, bo jego detektory
r\'ownie\.z skaluj\k{a} si\k{e} z~$\Phi$.
\end{remark}

\begin{remark}[Kosmologiczna zachowawczo\'s\'c entropii]%
\label{rem:entropy-conservation}
G\k{e}sto\'s\'c entropii promieniowania $s_{\mathrm{rad}}
= (2\pi^2/45)\,k_B\,(k_B T/(\hbar c))^3$ skaluje si\k{e} jako
$s_{\mathrm{rad}}(\Phi) = s_{\mathrm{rad},0}\,(\Phi/\PhiZero)^{3/2}$
przy ustalonym~$T$. W~kosmologii FRW z~jednorodnym $\Phi = \PhiZero\psi$
(gdzie $\psi$ \'sledzi atraktor --- por.~\S\ref{ssec:friedmann})
ekspansja adiabatyczna zachowuje $s\,a^3 = \mathrm{const}$,
bo $\psi$ jest jednorodne i~nie generuje gradient\'ow entropii.
Ewolucja kosmologiczna $\psi$ jest dostatecznie powolna,
by lokalna r\'ownowaga termodynamiczna by\l{}a utrzymana w~ka\.zdej epoce.

Kowariantna zachowawczo\'s\'c $\nabla_\mu T^{\mu\nu} = 0$
obowi\k{a}zuje dla materii sprz\k{e}\.zonej z~metryk\k{a} TGP
(aksjomat~\ref{ax:metric-coupling}). Pozorna wymiana energii
mi\k{e}dzy $\Phi$ a~materi\k{a} jest ujmowana przez r\'ownanie
polowe --- energia nie jest ,,tracona''.
\end{remark}

% ----------------------------------------------------------
\subsection{Otwarty problem}\label{ssec:stale-problem}

\begin{problem}[Obserwowalno\'s\'c dynamicznych sta\l{}ych]\label{prob:stale}
Zale\.zno\'sci $c(\Phi)$, $\hbar(\Phi)$, $G(\Phi)$ s\k{a} testowalne
tylko wtedy, gdy $\Phi$ zmienia si\k{e} mierzalnie mi\k{e}dzy r\'o\.znymi
regionami. W~warunkach laboratoryjnych $\Phi \approx \PhiZero$
wsz\k{e}dzie, wi\k{e}c odchylenia s\k{a} minimalne. Predykcje:
\begin{itemize}[nosep]
  \item efekty powinny by\'c mierzalne w~silnych polach
        (blisko gwiazd neutronowych, czarnych dziur),
  \item kosmologicznie: $\Phi$ w~p\'o\'znym Wszech\'swiecie
        mo\.ze r\'o\.zni\'c si\k{e} od $\Phi$ we wczesnym ---
        obserwowalne przez stosunek $G_{\mathrm{wczesny}}/G_{\mathrm{teraz}}$.
\end{itemize}
\end{problem}

% ----------------------------------------------------------
\subsection{Interpretacja w~kontek\'scie SI i~obserwabli bezwymiarowych}%
\label{ssec:si-metrologia}

\begin{remark}[Stale fundamentalne a~definicje SI]\label{rem:SI-constants}
W~uk\l{}adzie SI (BIPM, 2019) sta\l{}e $c$ i~$\hbar$ maj\k{a}
\emph{ustalone dok\l{}adne warto\'sci liczbowe}: $c = 299\,792\,458$~m/s,
$\hbar = 1{,}054\,571\,817\ldots \times 10^{-34}$~J$\cdot$s.
Twierdzenie o~,,zmiennym $c$'' nie oznacza zmiany tej liczby
--- oznacza zmian\k{e} \emph{relacji mi\k{e}dzy procesami fizycznymi}
mierzon\k{a} w~jednostkach SI.

\medskip\noindent
\textbf{Precyzyjnie:} aksjomat A6 ($c(\Phi) = c_0\sqrt{\PhiZero/\Phi}$)
m\'owi, \.ze \emph{metryka efektywna}, w~kt\'orej propaguje si\k{e}
\'swiat\l{}o, zale\.zy od pola $\Phi$. Obserwator lokalny
(u\.zywaj\k{a}cy zegar\'ow i~pr\k{e}t\'ow sprzegni\k{e}tych
z~tym samym $\Phi$) zmierzy \emph{zawsze} $c = c_0$
--- to jest wymaganie zasady r\'ownowa\.zno\'sci (EEP).

\medskip\noindent
\textbf{Obserwowalne s\k{a} tylko bezwymiarowe kombinacje:}
\begin{itemize}[nosep]
  \item $\alpha_{\mathrm{em}} = e^2/(4\pi\varepsilon_0 \hbar c)$
    --- je\'sli $\hbar$ i~$c$ zmieniaj\k{a} si\k{e} tak samo
    (aksj.~A6--A7), to $\alpha_{\mathrm{em}}$
    \emph{nie dryfuje} z~$\Phi$.
  \item $\mu \equiv m_p/m_e$ --- stosunek mas niezale\.zny od $\Phi$
    (masy skaluj\k{a} si\k{e} jak $m \propto A_{\rm tail}^4$,
    wsp\'olna skala $\Phi$-zale\.zna si\k{e} redukuje).
  \item $G(\Phi)\,m_p^2/(\hbar c)$ --- \emph{ten} stosunek dryfuje:
    $\propto \PhiZero/\Phi$. Jest to jedyna bezwymiarowa kombinacja
    w~TGP zmieniaj\k{a}ca si\k{e} z~$\Phi$.
\end{itemize}

\medskip\noindent
\textbf{Testowalno\'s\'c:}
dryf $G$ wzgl\k{e}dem ska\l{} atomowych mierzony jest przez
\emph{Lunar Laser Ranging} (LLR: $|\dot G/G| < 2 \times 10^{-2}\,H_0$)
i~obserwacje BBN ($\Delta G/G \lesssim 0{,}05$).
TGP daje $|\dot G/G|/H_0 = 0{,}009$ (numerycznie,
\texttt{p\_frw\_full\_evolution.py}), co jest \emph{zgodne}
z~obecnymi ograniczeniami.

\medskip\noindent
\textbf{Podsumowuj\k{a}c:} TGP \emph{nie} postuluje zmienno\'sci
sta\l{}ych fizycznych w~sensie zmiany definicji SI,
lecz zmienno\'s\'c \emph{metryki efektywnej} i~\emph{sprz\k{e}\.ze\'n},
prowadz\k{a}c\k{a} do jednej testowalnej konsekwencji:
dryfu $G\,m_p^2/(\hbar c)$ na poziomie $\lesssim 1\%$ na $H_0^{-1}$.
\end{remark}

% ----------------------------------------------------------
\subsection{Mechanizm hierarchii mas ferm\-ion\'ow}\label{ssec:yukawa-kink}
% v24: F-III-1, scripts/yukawa_mass_hierarchy_v24.py (15/15 PASS)

\begin{proposition}[Mechanizm wzmocnienia Yukawa]\label{prop:yukawa-kink}
Masa $m_n$ fermion\'ow $n$-tej generacji ($n=0,1,2$) sk\l{}ada si\k{e}
z~dw\'och czynnik\'ow:
\begin{equation}\label{eq:yukawa-kink}
  m_n \;=\; m_0 \cdot \frac{E_n}{E_0}
        \cdot \exp\!\bigl(\kappa_Y \cdot \Delta V_{\mathrm{kin}}^{(n)}\bigr),
\end{equation}
gdzie:
\begin{itemize}[nosep]
  \item $E_n$ --- energia WKB $n$-tego stanu zwi\k{a}zanego w~potencjale
    Sturm--Liouville'a kinku substratu
    ($E_0{:}E_1{:}E_2 \approx 1{:}2.14{:}2.79$, por.\
    dodatek~F, prop.~\ref{prop:no-4th-generation}),
  \item $\Delta V_{\mathrm{kin}}^{(n)}
         = V_{\mathrm{kin}}^{(n)} - V_{\mathrm{kin}}^{(0)}$
    --- przyrost obj\k{e}to\'sci kinetycznej kinku:
    \(
      V_{\mathrm{kin}}^{(n)}
      = \int_0^\infty \!\bigl(\tfrac{d\chi_n}{d\xi}\bigr)^2\xi^2\,d\xi,
    \)
  \item $\kappa_Y$ --- bezwymiarowy parametr sprzężenia substratu
    z~sektorem elektros\l{}abym, zdeterminowany przez
    $\{J,\,C_{\mathrm{kin}},\,C_\gamma\}$.
\end{itemize}
Z~obserwowanych mas lepton\'ow:
\begin{equation}\label{eq:kappaY}
  \kappa_Y = \frac{\ln(m_\mu/m_e)/(E_1/E_0)}{\Delta V_{\mathrm{kin}}^{(1)}}
           \approx 0.061\,V_{\mathrm{kin}}^{-1}
           \quad (\text{spójne z~}\kappa_Y(m_\tau)\ \text{do}\ 9\%),
\end{equation}
z~predykcją $m_\tau/m_e \approx 2.0 \times 3477 = 6954$
(dok\l{}adno\'s\'c rzędu wielkości; pełna predykcja wymaga MC).
\end{proposition}

\begin{remark}[Status O16 -- v24]\label{rem:O16-v24}
Mechanizm wzmocnienia Yukawa (prop.~\ref{prop:yukawa-kink}) jest
\textbf{ustalony} (v24): dwa sk\l{}adniki hierarchii $m_n$ zidentyfikowane,
$\kappa_Y$ wyznaczony z~obserwacji do precyzji~$9\%$. Problem O16
pozostaje \textbf{cz\k{e}\'sciowo otwarty}: precyzyjne wartości
$V_{\mathrm{kin}}^{(n)}$ dla $n=1,2$ wymagaj\k{a} BVP z~w\k{e}z\l{}ami
oraz symulacji FSS MC~3D ($N^3 \geq 32^3$).
Weryfikacja numeryczna: \texttt{scripts/yukawa\_mass\_hierarchy\_v24.py}
-- 15/15~PASS\@.
\end{remark}

% --- Dimensionless observables and metrological protocol (H5 deliverable) ---
% tgp_metrology_observables.tex -- Dimensionless observables and metrological protocol
% Self-contained section for \input{} into main manuscript.
% References: BIPM 2019 (SI redefinition), CODATA 2018/2022.

\subsection{Dimensionless observables and metrological consistency}%
\label{ssec:metrology}

\paragraph{The problem.}
TGP posits $c(\Phi)$, $\hbar(\Phi)$, $G(\Phi)$ as field-dependent.
Since the 2019 SI fixes $c$ and $h$ as \emph{exact constants}
(BIPM, Resolution~1 of the 26th CGPM), any claim that these
quantities ``vary'' must be stated carefully.

\paragraph{Resolution.}
The physical content of TGP is fully expressed through
\textbf{dimensionless combinations} that are convention-independent:

\begin{center}\small
\begin{tabular}{@{}llll@{}}
\toprule
\textbf{Observable} & \textbf{Expression} & \textbf{TGP dependence on $\Phi$} & \textbf{Test} \\
\midrule
Fine structure constant $\alpha$
  & $e^2/(4\pi\varepsilon_0\hbar c)$
  & $\alpha(\Phi) = \alpha_0$ \quad (\textbf{constant})
  & Atomic clocks; quasar absorption \\[4pt]
Proton-to-electron mass ratio $\mu$
  & $m_p/m_e$
  & $\mu(\Phi) = \mu_0$ \quad (\textbf{constant})
  & Molecular spectroscopy \\[4pt]
Gravitational fine structure $\alpha_G$
  & $G m_p^2/(\hbar c)$
  & $\alpha_G(\Phi) = \alpha_{G,0}\,/\psi^2$ \quad (\textbf{varies})
  & Binary pulsar timing \\[4pt]
Gravitational redshift $z_{\rm grav}$
  & $\nu_A/\nu_B - 1$
  & $z_{\rm grav} = \sqrt{\psi_B/\psi_A} - 1$
  & Gravity Probe A; ACES; GPS \\[4pt]
Hubble rate $H(z)$
  & $\dot{a}/a$
  & $H(z;\Phi_0,\gamma)$ (modified Friedmann)
  & BAO (DESI); CMB (Planck) \\[4pt]
GW/EM speed ratio
  & $c_{\rm GW}/c_{\rm EM}$
  & $= 1$ \quad (exact, algebraic identity)
  & Multi-messenger (GW170817) \\
\bottomrule
\end{tabular}
\end{center}

\paragraph{Key point.}
In TGP, $\alpha$ and $\mu$ are \textbf{$\Phi$-independent}
because both $e^2/(\hbar c)$ and $m_p/m_e$ involve ratios
where the $\Phi$-dependence cancels.  This is a non-trivial
consistency condition: TGP predicts that atomic physics is
locally invariant even though individual ``constants'' vary.
Only \emph{gravitational} dimensionless quantities ($\alpha_G$,
$z_{\rm grav}$) carry $\Phi$-dependence.

\paragraph{CODATA protocol.}
All numerical comparisons in this work use:
\begin{itemize}[nosep]
  \item CODATA 2018 recommended values for fundamental constants
        (NIST SP~961, 2019 edition);
  \item PDG 2024 for particle masses and coupling constants;
  \item Planck 2018 (VI, Table~2) for cosmological parameters
        ($H_0 = 67.36 \pm 0.54$\,km/s/Mpc, $\Omega_m = 0.3153$, etc.);
  \item DESI DR1 BAO distances for $H(z)$ and $D_M(z)$ constraints
        (DESI Collaboration 2024, arXiv:2404.03002).
\end{itemize}
When TGP and GR predictions differ, the comparison is always
stated as a deviation in a \textbf{dimensionless observable}
(not in a dimensional ``constant'').

\paragraph{Experimental sensitivity.}
Current bounds on variation of $\alpha$ and $\mu$:
$|\dot\alpha/\alpha| < 10^{-17}\,{\rm yr}^{-1}$ (optical clocks, PTB/NIST),
$|\dot\mu/\mu| < 10^{-16}\,{\rm yr}^{-1}$ (molecular spectroscopy).
TGP predicts $\dot\alpha = \dot\mu = 0$ at all $\Phi$ ---
these bounds are \textbf{automatically satisfied}.
The non-trivial TGP signal is in $\alpha_G(\Phi)$: the effective
gravitational coupling varies as $G_{\rm eff} \propto 1/\psi^2$.
For $\psi = 1 + 2U/c^2$, the deviation is
$\delta G/G_0 \approx -4U/c^2 \sim -6 \times 10^{-10}$
on the surface of the Sun --- consistent with helioseismology
bounds ($|\delta G/G| < 10^{-4}$, Guenther et al.\ 1998).


% --- Time glossary and EEP analysis (H1 deliverable) ---
% tgp_time_glossary.tex -- Glossary of time concepts and EEP analysis
% Self-contained section for \input{} into main manuscript.
% Synchronized with sek08c, sek08 (rem:substrate-time), sek04 (ax:c).

\subsection{Glossary: substrate time, proper time, and clocks}%
\label{ssec:time-glossary}

TGP introduces a non-standard time structure.  To prevent confusion,
we define all temporal quantities explicitly.

\begin{center}\small
\begin{tabular}{@{}lp{4.5cm}p{6.5cm}@{}}
\toprule
\textbf{Symbol} & \textbf{Name} & \textbf{Definition and role} \\
\midrule
$dt_{\rm sub}$
  & Substrate tick
  & Global, $\Phi$-independent increment of the discrete substrate clock.
    Not directly measurable; serves as the ``absolute time'' of the lattice
    $\Gamma$.  All observers count the same number of substrate ticks
    between two events. \\[6pt]
$d\tau$
  & Proper time
  & Time measured by a physical clock (e.g.\ Cs hyperfine transition).
    Depends on local $\Phi$:
    \begin{equation*}
      d\tau = \frac{c_0}{\sqrt{\psi}}\,dt_{\rm sub},
      \qquad \psi \equiv \Phi/\Phi_0.
    \end{equation*}
    In regions with $\Phi > \Phi_0$ (``denser space''): $d\tau < c_0\,dt_{\rm sub}$
    --- clocks run slower. \\[6pt]
$c_{\rm lok}$
  & Local speed of light
  & Measured by a \emph{local} observer using their proper time and
    proper length:
    $c_{\rm lok} = c_0/\sqrt{\psi}$.
    This is \textbf{not} the same as the coordinate speed
    $c_{\rm coord} = c_0/\psi$ (which is frame-dependent). \\[6pt]
$c_0$
  & Reference speed
  & Speed of light at $\Phi = \Phi_0$ (reference vacuum).
    Identified with the SI value $c = 299\,792\,458$\,m/s.
    In TGP, $c_0$ is a \emph{property of the reference state},
    not a universal constant. \\[6pt]
$\nu_{\rm Cs}(\Phi)$
  & Cs clock frequency
  & Hyperfine transition frequency of $^{133}$Cs (SI second definition).
    In TGP: $\nu_{\rm Cs}(\Phi) = \nu_{\rm Cs,0}\,\sqrt{\Phi_0/\Phi}$.
    Two clocks at different $\Phi$ tick at different rates
    --- this is the gravitational redshift, here derived from substrate
    budget rather than postulated. \\[6pt]
$T_{\rm period}$
  & Orbital/dynamical period
  & Period of a macroscopic orbit (e.g.\ planet).
    Measured in substrate ticks, it is $\Phi$-independent.
    Measured in proper time, it depends on $\Phi$. \\
\bottomrule
\end{tabular}
\end{center}

\paragraph{Einstein Equivalence Principle (EEP).}
TGP maintains a \emph{modified} form of EEP:
\begin{enumerate}[nosep]
  \item \textbf{Weak Equivalence Principle (WEP):}
        \textsc{Satisfied.}  In TGP the inertial mass equals the
        gravitational source strength ($m_{\rm inert} = C_i = m_{\rm grav}$);
        all bodies fall identically in a given $\nabla\Phi$.
  \item \textbf{Local Lorentz Invariance (LLI):}
        \textsc{Satisfied locally.}  In a sufficiently small region,
        $\Phi \approx \text{const}$ and the metric reduces to Minkowski
        with $c = c_{\rm lok}$.  No preferred direction arises.
  \item \textbf{Local Position Invariance (LPI):}
        \textsc{Modified.}  Non-gravitational experiments
        (e.g.\ fine structure constant $\alpha$, mass ratios)
        are $\Phi$-independent in TGP --- they depend on
        dimensionless combinations only.  But $c$, $\hbar$, $G$
        individually depend on $\Phi$.
        The \emph{observable} consequence: gravitational redshift
        $z_{\rm grav} = \sqrt{\Phi_{\rm emit}/\Phi_{\rm obs}} - 1$
        follows the \textbf{same} formula as GR (to leading order
        in $\delta\Phi/\Phi_0$).
\end{enumerate}

\paragraph{Operational protocol.}
To compare TGP predictions with clock experiments
(GPS, Gravity Probe A, ACES, MICROSCOPE):
\begin{enumerate}[nosep]
  \item Express all observables as \textbf{dimensionless frequency ratios}:
        $\nu_A/\nu_B$ between clocks A and B.
  \item In TGP: $\nu_A/\nu_B = \sqrt{\Phi_B/\Phi_A} = \sqrt{\psi_B/\psi_A}$.
  \item In GR: $\nu_A/\nu_B = \sqrt{g_{tt}^{(B)}/g_{tt}^{(A)}}
        \approx 1 + (U_A - U_B)/c^2$.
  \item To leading order in $\delta\Phi/\Phi_0$: \textbf{identical}.
        Deviations appear at $O(\delta\Phi^2/\Phi_0^2)$, i.e.\
        $O(U^2/c^4) \sim 10^{-18}$ on Earth --- currently below
        experimental sensitivity.
\end{enumerate}

\paragraph{Relation to SI.}
In the 2019 SI, the values of $c$ and $h$ are \emph{exact by definition}:
$c = 299\,792\,458$\,m/s, $h = 6.626\,070\,15 \times 10^{-34}$\,J\,s.
TGP does \emph{not} contradict this: the SI definitions fix the
\emph{numerical value} of $c$ in the unit system, regardless of whether
$c$ is a universal constant or a field.  The physical content of TGP
is encoded in \textbf{dimensionless observables} (frequency ratios,
fine structure constant, mass ratios) that are measurement-convention-independent.


